\frfilename{mt12.tex}
\versiondate{7.4.05}
\copyrightdate{1994}

\def\chaptername{Integrasi}

\newchapter{12}

Jika Anda menyusuri rak yang sesuai di perpustakaan perguruan tinggi Anda,
Anda akan melihat bahwa kata `ukuran' dan `integrasi' berdampingan seperti
kembar siam. Keterkaitannya lebih rumit sekaligus lebih erat daripada yang
dapat digambarkan oleh penjelasan sederhana mana pun. Namun, jika kita
mengatakan bahwa salah satu konsep yang mendasari integrasi ialah `luas di
bawah kurva', jelas bahwa setiap metode untuk menentukan `luas' semestinya
bersesuaian dengan suatu metode untuk mengintegralkan fungsi; dan sejak
awal hal ini telah menjadi bagian penting dari teori Lebesgue. Untuk uraian
harfiah mengenai integral fungsi nonnegatif dalam kaitannya dengan luas
himpunan ordinatnya, menurut saya sebaiknya kita menunggu hingga Bab 25
dalam Jilid 2. Dalam bab ini saya berusaha memberikan uraian ringkas
mengenai integral baku suatu fungsi bernilai real pada ruang ukur umum,
beserta enam teorema terpenting tentang integral ini.

Konstruksi ini dipenuhi kesulitan teknis pada setiap langkah, dan Anda akan
mudah memahami mengapa konstruksi tersebut belum dilakukan sebelum tahun
1901. Hal yang mungkin kurang jelas ialah mengapa konstruksi itu akhirnya
dilakukan juga. Karena itu, barangkali Anda sebaiknya segera membaca
pernyataan 123A-123D %123A 123B 123C 123D
di bawah ini. Memang benar (sebagian perinciannya akan muncul jauh kemudian,
dalam \S436 di Jilid 4) bahwa setiap teori integrasi yang cukup kuat untuk
memiliki teorema-teorema semacam ini pada hakikatnya harus mencakup semua
gagasan dalam bab ini dan hampir semua gagasan dalam bab sebelumnya.

\discrpage
